% Generated draft: requires mathematical and source review.
\documentclass[11pt]{article}
\usepackage[margin=25mm]{geometry}
\usepackage{fontspec}
\setmainfont{Latin Modern Roman}
\usepackage{amsmath,amssymb,mathtools,unicode-math}
\setmathfont{Latin Modern Math}
\usepackage{enumitem,needspace,xurl,hyperref}
\setlength{\parindent}{0pt}
\setlength{\parskip}{5pt}
\setlength{\emergencystretch}{3em}
\newcommand{\N}{\mathbb{N}}
\newcommand{\Z}{\mathbb{Z}}
\newcommand{\Q}{\mathbb{Q}}
\newcommand{\R}{\mathbb{R}}
\newcommand{\C}{\mathbb{C}}
\newcommand{\F}{\mathbb{F}}
\newcommand{\A}{\mathbb{A}}
\newcommand{\PP}{\mathbb{P}}
\newcommand{\E}{\mathbb{E}}
\newcommand{\eps}{\varepsilon}
\newcommand{\dd}{\,\mathrm{d}}
\begin{document}
\section*{UnsolvedMath ID: 11000150}
\textbf{Source problem code:} catalogue entry
\subsection{Definitions and mathematical statement}
Let $S$ be an oriented genus-$g$ surface with exactly one boundary component $\delta$, and use the boundary-pointwise mapping class group. Let $T_\alpha$ denote a positive Dehn twist. Smith's question asks for the maximal length $n$ of a factorization
\[
 T_{\alpha_1}\cdots T_{\alpha_n}=T_\delta
\]
whose curve configuration has pairwise intersection $0$ or $1$. In the usual minimal-position interpretation this means $i(\alpha_i,\alpha_j)\in\{0,1\}$ for every $i\ne j$, where $i$ is geometric intersection number. Repeated curves are allowed, explicitly so in the genus-two discussion.

After fixing the source's admissible curve class, define
\[
 L_g=\sup\{n\geq1:\text{such a length-}n\text{ factorization exists}\}.
\]
The problem asks to determine $L_g$ and whether arbitrarily long such products exist. The fixed-genus versus varying-genus scope should be confirmed. Essentiality and exclusion of boundary-parallel curves are not explicit in the excerpt; inessential identity twists would make unrestricted length artificially unbounded, so the intended restriction is essential to the question.

For genus $2$, the source says the curves belong to a $D_6$ configuration with repetitions and that its Coxeter element gives the square of a boundary twist. It tentatively suggests excluding the pictured curve $\alpha'_1$ and reducing to an $A_5$ configuration, then announces an algebraic question. Neither the configurations, the curve diagram nor the resulting algebraic question is supplied. These are retained as source suggestions, not an established reduction or an invented presentation.

The registry reports arbitrarily long positive boundary-multitwist factorizations in work of Baykur--Van Horn-Morris, while explicitly distinguishing the present one-boundary and pairwise-intersection restriction.
\subsection{Short English statement}
How long can a positive factorization of one boundary twist be if every pair of factor curves intersects at most once? Preserve the restriction and the unresolved genus-two configuration reduction.
\subsection{Sources}
\begin{itemize}
\item Local registry: \url{data/unsolvedmath/problems.json}, permanent UnsolvedMath ID $11000150$, exact source code \texttt{catalogue entry}.
\item Farb, \emph{Problems on Mapping Class Groups and Related Topics}, Question (Smith), PDF page 133: \url{https://www.math.uchicago.edu/~farb/papers/mcgbook.pdf}. Supplied access date: 2026-07-29; extraction: pdf-font-structure; heading disambiguated by page among 329 headings.
\item The registry attributes unrestricted long-factorization progress to Baykur--Van Horn-Morris, without precise references. The source's $D_6/A_5$ diagrams and announced algebraic question are absent.
\end{itemize}
\end{document}
